\scanpage{29}
\uuline{Lemme} Notations $\unsure{\mathbb T|T},F,n,\mathcal M$ comme dans le
lemme de Cartan.\inserted{Considérons les} Conditions \deleted{équivalentes}

(i) $\forall S\in\mathcal M$,
\figure{\begin{tikzpicture}[baseline=(n.base)]\node[draw,ellipse,inner sep=5pt] (n) {$\deleted{\unsure{H}}\varinjlim_{\mathcal U}H^i(\mathcal U,F)$};\end{tikzpicture}}
$\simeq H^i(S,F).\qquad i\leq n$

[$\mathcal U$ parcourt les rec. de $S$ formés de $S_i\in\mathcal M$,]

(ii) $\forall S\in\mathcal M$, \deleted{et tt rec. $\mathcal U$ de $S$ formé
d'objets de $\unsure{S|\mathcal M}$, et tt élément:}
\[
\varinjlim_{\mathcal U}H^p(\mathcal U;H^q(F))=0\quad\text{pour }
\deleted{\unsure{0<q\leq n|0<p+q\leq n}\text{ et }q\ne0}\inserted{0<q\leq n\unsure{+1}}
\]

(iii)
\[
\varinjlim_{\mathcal U}H^p(\mathcal U,\overline H^q(F))=0\quad\text{pour }0<q\leq n\unsure{+1}
\quad\deleted{\text{et }\unsure{p|q}=0}
\]
où on pose\inserted{pour $T\in\mathcal M$,}
\[
\overline H^q(F)(T)=\varinjlim_{\mathcal V}H^q(\mathcal V,T),
\]
$\mathcal V$ parcourant les rec. de $T$ formés d'objets de $\mathcal M$.

(ii bis) $\forall S\in\mathcal M$, $\mathcal U$ rec. de $S$ par des $U_i\in\mathcal M$,
et $\unsure{\gamma|\zeta}\in\unsure{C^p|Z^p}(\mathcal U,H^q(F))$,
\deleted{($\unsure{p\geq0},p+q\leq n-1,\ q\ne0$)}\inserted{$0<q\leq n$}
il existe $\mathcal U'$ rec. de $S$ par des $U_i\in\mathcal M$, majorant $\mathcal U$,
\deleted{\unsure{tel que le}} et $\mathcal U'\unsure{-|\longrightarrow}\mathcal U$, tel que l'image de
$\unsure{\gamma|\zeta}$ dans $C^p(\mathcal U',H^q(F))$ soit nulle.

(iii bis) \deleted{$\unsure{\Leftarrow|\Leftrightarrow}$} Comme (ii bis), mais $H^qF$,
remplacé par $\overline H^q(F)$.

Mais on a
\diagram{\begin{tikzcd}[column sep=large,row sep=small]
\deleted{\text{(ii bis)}}\arrow[r,Rightarrow]\arrow[d,Rightarrow]&\deleted{\text{(ii)}}\arrow[d,Rightarrow]\\
{}\arrow[r,dash]&{}
\end{tikzcd}}
\diagram{\begin{tikzcd}[column sep=large,row sep=small]
\text{(ii bis)}\arrow[r,Leftrightarrow]\arrow[d,Rightarrow]&\text{(iii bis)}\\
\text{(ii \deleted{\unsure{bis}})}\arrow[r,Leftrightarrow]\arrow[d,Rightarrow]&\text{(iii)}\\
\text{(i)}&
\end{tikzcd}}
\margin{\figure{\begin{tikzpicture}[x=1cm,y=1cm]
\draw (0,4.8)--(.02,4.1);
\draw[thin] (.15,3.7)..controls(-.3,3.6)and(-.2,3.5)..(-.16,3.1)--(-.04,2.2)..controls(.1,2.2)and(.3,2.25)..(.48,2.3);
\draw[thin] (.1,1.8)..controls(-.1,1.8)and(-.08,1.7)..(-.1,1.5)--(-.1,.1)..controls(-.08,0)and(.06,0)..(.2,0);
\end{tikzpicture}}}
\mspage{28}

